% relationship to the student
% academic and intellectual ability
% professional achievements
% academic potential for graduate work
% address sustainability for the program
% ability to commit to the rigorous curriculum
% emotional development

% In all cases, try to give some examples instead of just saying vague nice things.

% First paragraph (Intro)
% Introduction saying I recommend [full name] to [program] at [university].
% How long and in what capacity I know the student (taught in class, advised, was my TA, chaired thesis, etc.)
% Quick summary of attributes I think qualify the student.

I am pleased to give my strong recommendation for Hyeonbeen (Edward) Lee for admission to the \Program{} Ph.D. program at \University{}. With five years of close supervision and mentorship since 2021, I witnessed his entire academic journey---from a student in my class to an undergraduate researcher, a Master's student, and a postgraduate researcher. Through this prolonged interaction, I have developed solid confidence that Lee possesses outstanding intellectual curiosity, technical skills, and resilience required to excel in your program and as an independent researcher.

% Next section (classes they took with me)
% Brief description of course (survey, seminar, etc.)
% Overall assessment of their performance (not just the grade)
% Some examples of why they did well (some from an assignment, or exam performance, or class discussion ability, etc.)
% Repeat if more than one class

In my courses on computational dynamics and finite element methods, he was one of the finest students and consistently achieved top grades by excelling at understanding complex formulations and their computational implementations. He was especially impressive for his proactive analyses beyond the required tasks, such as using symbolic math libraries to demonstrate his derivation, or translating MATLAB codes into Python and sharing them via GitHub for reproducibility, all of which went well beyond the course material. These self-motivated activities, driven by his intellectual curiosity, strongly validated his potential as a researcher, which was my key reason to accept him as an undergraduate researcher.

% Next section (research experience)
% What projects the student was involved in
% Specifically what parts the student was involved in
% Quality of their work and level of supervision needed
% What are the outcomes of the research (presentations, publications, etc.)
% If the student presented work, I will describe the type venue and assess how they did

After joining the lab, he showed an intense interest in \textit{machine learning} and \textit{neural networks}. He was particularly passionate and proficient in our research on data-driven simulation. As an undergraduate researcher, he rapidly digested relevant publications and reproduced key results within months, while also leading a machine learning thesis project with his undergraduate co-advisor Prof. Shinkyu Jeong.

He always sought to produce something beyond the given task, which distinguished him from other students. While co-working with our team on the vehicle dynamics simulation and neural network-based meta-modeling, 
% which was a project with the Korea Atomic Energy Research Institute (KAERI),
he identified the limitation of our existing approach in modeling impulsive dynamics and proposed a novel multi-order learning architecture. This work eventually led to his first-author publication in \textit{Journal of Computational Physics}, which also received  the Excellent Paper Award from the Korean Society of Mechanical Engineers conference.
% he noticed that networks struggled to learn impulsive high-order dynamics, such as the vertical acceleration of a vehicle passing a bump. To address this, he independently studied the spectral bias of neural networks and suggested simultaneous learning of smooth lower-order dynamics and volatile higher-order dynamics
% at the end of his undergraduate studies. This idea formed the basis for his first-author publication `cNN-DP', which was not only published in the prestigious \textit{Journal of Computational Physics} but also received an Excellent Paper Award from the Korean Society of Mechanical Engineers (KSME) conference.

At the end of his first Master's year, he also started working on the project of ML-based joint force estimation for a marine excavation robot. 
% with the Korea Research Institute of Ships and Ocean Engineering (KRISO). 
Our initial objective was a single-step force prediction in a simulated environment. However, within a year of him taking on the project,
Lee progressively expanded the project by shifting to real-world sensor data, applying sequential modeling, and reframing it as a forecasting problem. Consequently, he came up with a decomposition-based generative forecasting framework for contact-rich real-world force dynamics. 
% he significantly expanded this work by (1) switching the learning domain to the real world, (2) introducing sequence-based modeling, (3) redefining it as a multi-horizon forecasting problem, and (4) developing a decomposition-based generative strategy for forecasting high-frequency, contact-rich robot joint forces. 
These improvements yielded not only an impressive advance in the model's accuracy but also served as a valuable foundation for the lab's machine learning capability and led to subsequent success in the project.

% In his final Master's semester, his interests further expanded to \textit{reinforcement learning} and \textit{robotics}, and he earned an opportunity to work with one of the nation's best groups. However, he declined the offer, believing it was not the right time, and instead focused more on his thesis and revisions. 
After earning his Master's degree, Lee continued on this project as a postgraduate researcher.
For two years, he dedicated his effort to accelerate the model's inference for real-time application through systematic optimization, along with in-depth refinement of the model by testing and integrating state-of-the-art learning methodologies. 
% This refinement involved experiments such as leveraging self-supervised learning, the physical interpretation of learned parameters, positional placement of layers, validation of state-of-the-art model backbones, and more. 
Concurrently, he also enthusiastically mentored a junior master's student who was building upon his work.
% Based on these long-term research experience with him, I can firmly attest to his perseverance, responsibility, and deep-seated passion for doctoral studies.

% For two years, he worked closely with our group on (1) acceleration of inference for real time application by architectural compression, compiling and parallelization, (2) in-depth refinement of the model by experiments such as leveraging self-supervised learning, physical interpretation of learned parameters, positional placement of layers, validation of state-of-the-art model backbones, and more. Meanwhile, he also mentored a junior master's student who was building on his work. Based on these experiences, I can firmly attest to his perseverance, responsibility, and deep-seated passion for doctoral studies.

% Next section (general attributes)
% Personal strengths
% How well they work with others
% How well they take instruction / feedback
% Comment on the students future plans and if they are moving toward them
% Comment on hardships they have overcome (work full time, first in family in college, etc.) or address a red flag if needed (why they did poor only one semester, etc.)
% Few more personal strengths (always end on positive stuff)
% Compare students to other students I have taught
% Give bottom line assessment where I again recommend the person for the program

Besides his research excellence, he also maintained good relationships with the other lab members and shared his knowledge with the team. He created the lab's GitHub, led systematic code management, and actively shared well-documented repositories to support his peers. Furthermore, he used to get along well with international and visiting lab members and helped them overcome environmental hardships, all of which can demonstrate his emotional maturity and team spirit.

In summary, Lee possesses an outstanding combination of intellectual passion, validated research skills, and collaborative spirit. He is one of the most impressive students I have mentored, and I am confident he will make an excellent contribution to your Ph.D. program. Therefore, I recommend him without any reservation.